\subsection{The odd minimum is chosen after its lower bound}
\begin{proposition}[The odd minimum and every finite operator budget]\label{ex:app-odd-minimum}
Fix the original odd filter, the detecting sequence of
Proposition~\ref{ex:app-detector}, and one nonprincipal ultrafilter
on its stages. The arithmetic lower bound (3.5), the native comparisons
(3.4), and Proposition~\ref{ex:app-rough-companion} give (F.4)--(F.5)
before any odd minimum is chosen. The bounded-weight visible depths
therefore have an attained finite integer minimum $\delta_o$.
For every fixed weight cutoff, the simultaneous inequality (F.6)
holds on a set in this same ultrafilter. For every finite operator
budget and $R\ge2$, one may choose a stage where all the integer
divisions and residual errors required in
Proposition~\ref{ex:app-coefficient-rules} are valid at every Fourier
index.
On a fixed squarefree-output range the unit symbol defined before
(F.7) is exhaustive, and its value $1$ implies (F.7), with the one
constant $C_o=\delta_o+C_{\rm det}$. No minimum is changed on
omitting a vertex or adding stationary tests, and no single finite
stage is claimed to serve every unbounded weight range.
\end{proposition}
\begin{proof}\leavevmode
The arithmetic result (3.5) and the two native-lattice height
comparisons (3.4) give, using (F.3),
\[
 2j_{\rm gen}(h)\ge2\omega(h)+s(h)-C_{\rm gen},\qquad
 C_{\rm gen}=2A+C_R+C_E-v_2L_*(h_*)+C_H.                         \tag{F.4}
\]
The variable $s(hk')$ cancels. This estimate precedes the theta
prime, precision, odd minimum, final field and binary dimension.

Combining (F.4) with (O.8) gives on all admissible $D$
\[
            v_2a_M(D)\ge\min\{M,\omega(D)-C_{\rm coarse}\},
 \qquad C_{\rm coarse}=C_{\rm gen}/2+C_{\rm det}.                 \tag{F.5}
\]
At visible coefficients (O.8) first gives a simple zero, allowing
(F.4); at the other odd-sign parameters the coefficient is zero
modulo $2^M$. This avoids a circular use of the future minimum.

